\section{AMC 12 12B 2002}
        \begin{enumerate}
        	\item (\textit{AMC 12 12B 2002}) The arithmetic mean of the nine numbers in the set $\{9, 99, 999, 9999, \ldots, 999999999\}$ is a $9$-digit number $M$, all of whose digits are distinct. The number $M$ does not contain the digit


 $\mathrm{(A)}\ 0 \qquad\mathrm{(B)}\ 2 \qquad\mathrm{(C)}\ 4 \qquad\mathrm{(D)}\ 6 \qquad\mathrm{(E)}\ 8$

 \textbf{Answer: }A



\textbf{Solution 1}

We wish to find $\frac{9+99+\cdots +999999999}{9}$, or $\frac{9(1+11+111+\cdots +111111111)}{9}=123456789$. This doesn't have the digit 0, so the answer is $\boxed{\mathrm{(A)}\ 0}$




\textbf{Solution 2}

Notice that the final number is guaranteed to have the digits $\{1, 3, 5, 7, 9\}$ and that each of these digits can be paired with an even number adding up to 9. $\boxed{\mathrm{(A)}\ 0}$ can be taken out, with the other digits fulfilling divisibility by 9.

 



\textbf{Solution 3}

The arithmetic mean is $\frac{(10^1-1)+(10^2-1)+\ldots+(10^9-1)}{9}=\frac{1111111101}{9}=123456789$. So select $\boxed{\mathrm{A}}$.~hastapasta

 


	\item (\textit{AMC 12 12B 2002}) What is the value of 

 \[(3x - 2)(4x + 1) - (3x - 2)4x + 1\]


 when $x=4$?


 $\mathrm{(A)}\ 0 \qquad\mathrm{(B)}\ 1 \qquad\mathrm{(C)}\ 10 \qquad\mathrm{(D)}\ 11 \qquad\mathrm{(E)}\ 12$

 \textbf{Answer: }D



\textbf{Solution 1}

By the distributive property, 

 
 \[(3x-2)[(4x+1)-4x] + 1 = 3x-2 + 1 = 3x-1 = 3(4) - 1 = \boxed{\mathrm{(D)}\ 11}\].

 



\textbf{Solution 2}

Inputting 4 into $x$ in the original equation,

 
 \[[3(4)-2][4(4)+1]-[3(4)-2][4(4)]+1 = 170-160+1 = \boxed{\mathrm{(D)}\ 11}\]



	\item (\textit{AMC 12 12B 2002}) For how many positive integers $n$ is $n^2 - 3n + 2$ a prime number?


 $\mathrm{(A)}\ \text{none} \qquad\mathrm{(B)}\ \text{one} \qquad\mathrm{(C)}\ \text{two} \qquad\mathrm{(D)}\ \text{more\ than\ two,\ but\ finitely\ many} \qquad\mathrm{(E)}\ \text{infinitely\ many}$

 \textbf{Answer: }B



\textbf{Solution 1}

Factoring, we get $n^2 - 3n + 2 = (n-2)(n-1)$. Either $n-1$ or $n-2$ is odd, and the other is even.  Their product must yield an even number.  The only prime that is even is $2$, which is when $n$ is $3$ or $0$. Since $0$ is not a positive number, the answer is $\boxed{\mathrm{(B)}\ \text{one}}$.

 



\textbf{Solution 2}

Considering parity, we see that $n^2 - 3n + 2$ is always even.  The only even prime is $2$, and so $n^2-3n=0$ whence $n=3\Rightarrow\boxed{\mathrm{(B)}\ \text{one}}$.

 


	\item (\textit{AMC 12 12B 2002}) Let $n$ be a positive integer such that $\frac 12 + \frac 13 + \frac 17 + \frac 1n$ is an integer. Which of the following statements is \textbf{not} true:


 $\mathrm{(A)}\ 2\ \text{divides\ }n \qquad\mathrm{(B)}\ 3\ \text{divides\ }n \qquad\mathrm{(C)}\ 6\ \text{divides\ }n \qquad\mathrm{(D)}\ 7\ \text{divides\ }n \qquad\mathrm{(E)}\ {n > 84}$

 \textbf{Answer: }E



\textbf{Solution 1}

Since $\frac 12 + \frac 13 + \frac 17  = \frac {41}{42}$, $0 < \lim_{n \rightarrow \infty} \left(\frac{41}{42} + \frac{1}{n}\right) < \frac {41}{42} + \frac 1n < \frac{41}{42} + \frac 11 < 2$

 From which it follows that $\frac{41}{42} + \frac 1n = 1$ and $n = 42$. The only answer choice that is not true is $\boxed{\mathrm{(E)}\ n>84}$.

 



\textbf{Solution 2 (no limits)}

Since $\frac 12 + \frac 13 + \frac 17  = \frac {41}{42}$, it is very clear that $n=42$ makes the expression an integer*. Because $n$ is a positive integer, $\frac{1}{n}$ must be less than or equal to $1$. Thus, the only integer the expression can take is $1$, which means the only value for $n$ is $42$. Thus $\boxed{\mathrm{(E)}\ n>84}$




\textbf{Solution 2.1 (faster ending to solution 2)}

Once you find $n=42$, you can skip the rest of Solution 2 by noting that $n=42=2\times3\times7$, which satisfies A, B, C, D, but not E. Thus $\boxed{\mathrm{(E)}\ n>84}$ must be the answer by PoE (Process of Elimination).

 



\textbf{Solution 3(similiar to solution2)}

Cross multiplying and adding the fraction we get the fraction to be equal to $\frac {41n + 42}{42n}$, This value has to be an integer. This implies,

 $42n|(41n+42)$.

 => $42|(41n + 42)$
\begin{verbatim}
but $42|42$, hence $42|41n$  but $42$ does not divides $41$,  $42|n$                                                                                                    -(1)
\end{verbatim}

=> $n|(41n+42)$
\begin{verbatim}
but $n|41n$,  $n|42$                                                                                                    -(2)
\end{verbatim}

from (1) and (2) we get that n=42.Comparing this with the options, we see that option E is the incorrect statement and hence E is the answer.

 


	\item (\textit{AMC 12 12B 2002}) Let $v, w, x, y,$ and $z$ be the degree measures of the five angles of a pentagon. Suppose that $v < w < x < y < z$ and $v, w, x, y,$ and $z$ form an arithmetic sequence. Find the value of $x$.


 $\mathrm{(A)}\ 72 \qquad\mathrm{(B)}\ 84 \qquad\mathrm{(C)}\ 90 \qquad\mathrm{(D)}\ 108 \qquad\mathrm{(E)}\ 120$

 \textbf{Answer: }D



\textbf{Solution 1}

The sum of the degree measures of the angles of a pentagon (as a pentagon can be split into $5- 2 = 3$ triangles) is $3 \cdot 180 = 540^{\circ}$. If we let $v = x - 2d, w = x - d, y = x + d, z = x+2d$, it follows that 

 
 \[(x-2d)+(x-d)+x+(x+d)+(x+2d) = 5x = 540 \Longrightarrow x = 108 \ \mathrm{(D)}\]

 Note that since $x$ is the middle term of an arithmetic sequence with an odd number of terms, it is simply the average of the sequence.

 You can always assume the values are the same so $\frac{540}{5}=108$




\textbf{Solution 2}

Let $v$, $w$, $x$, $y$, $z$ be $v$, $v + d$, $v+2d$, $v+3d$, $v+4d$, respectively. Then we have 
 \[v + w + x + y + z = 5v + 10d = 180^{\circ} (5 - 2) = 540^{\circ}\]Dividing the equation by $5$, we have 
 \[v + 2d = x = 108^{\circ} \mathrm {(D)}\]



	\item (\textit{AMC 12 12B 2002}) Suppose that $a$ and $b$ are nonzero real numbers, and that the equation $x^2 + ax + b = 0$ has solutions $a$ and $b$. Then the pair $(a,b)$ is


 $\mathrm{(A)}\ (-2,1) \qquad\mathrm{(B)}\ (-1,2) \qquad\mathrm{(C)}\ (1,-2) \qquad\mathrm{(D)}\ (2,-1) \qquad\mathrm{(E)}\ (4,4)$

 \textbf{Answer: }C



\textbf{Solution 1}

Since $(x-a)(x-b) = x^2 - (a+b)x + ab = x^2 + ax + b = 0$, it follows by comparing \href{https://artofproblemsolving.com/wiki/index.php?title=Coefficient}{coefficients} that $-a - b = a$ and that $ab = b$. Since $b$ is nonzero, $a = 1$, and $-1 - b = 1 \Longrightarrow b = -2$. Thus $(a,b) = \boxed{\mathrm{(C)}\ (1,-2)}$.

 



\textbf{Solution 2}

Another method is to use \href{https://artofproblemsolving.com/wiki/index.php?title=Vieta\%27s\_formulas}{Vieta's formulas}. The sum of the solutions to this polynomial is equal to the opposite of the $x$ coefficient, since the leading coefficient is 1; in other words, $a + b = -a$ and the product of the solutions is equal to the constant term (i.e, $a*b = b$). Since $b$ is nonzero, it follows that $a = 1$ and therefore (from the first equation), $b = -2a = -2$. Hence, $(a,b) = \boxed{\mathrm{(C)}\ (1,-2)}$




\textbf{Solution 3}

Since $a$ and $b$ are solutions to the given equation, we can write the two equations $a^2 + a^2 + b = 2a^2 + b = 0,$ and $b^2 + ab + b = 0.$ From the first equation, we get that $b = -2a^2,$ and substituting this in our second equation, we get that $4a^4 - 2a^3 - 2a^2 = 0,$ and solving this gives us the solutions $a = -\frac{1}{2},$ $a = 0$ and $a = 1.$ We discard the first two solutions, as the first one doesnt show up in the answer choices and we are given that $a$ is nonzero. The only valid solution is then $a = 1,$ which gives us $b = -2,$ and $(a, b) = \boxed{\text{(C) } (1, -2)}$




\textbf{Solution 4 (Using the Answer Choices)}

Note that for roots $a$ and $b$, $ab = b$. This implies that $a=1$, and there is only one answer choice with $1$ in the position for $a$, hence, $(a,b) = \boxed{\mathrm{(C)}\ (1,-2)}$




\textbf{Solution 5(Answer Choices)}

We can use Vieta's formulas to solve this. We know that $a + b = -a$, and $ab = b.$ We can test the answer choices, so that we find that $\boxed{\mathrm{(C)}\ (1,-2)}$ works.

 -monkey\_land

 


	\item (\textit{AMC 12 12B 2002}) The product of three consecutive positive integers is $8$ times their sum. What is the sum of their squares?


 $\mathrm{(A)}\ 50 \qquad\mathrm{(B)}\ 77 \qquad\mathrm{(C)}\ 110 \qquad\mathrm{(D)}\ 149 \qquad\mathrm{(E)}\ 194$

 \textbf{Answer: }B



\textbf{Solution 1}

Let the three consecutive positive integers be $a-1$, $a$, and $a+1$. Since the mean is $a$, the sum of the integers is $3a$. So $8$ times the sum is just $24a$. With this, we now know that $a(a-1)(a+1)=24a\Rightarrow(a-1)(a+1)=24$.  $24=4\times6$, so $a=5$. Hence, the sum of the squares is $4^2+5^2+6^2=\boxed{\mathrm{ (B)}\ 77}$.

 


	\item (\textit{AMC 12 12B 2002}) Suppose July of year $N$ has five Mondays. Which of the following must occur five times in August of year $N$? (Note: Both months have 31 days.)


 $\mathrm{(A)}\ \text{Monday} \qquad\mathrm{(B)}\ \text{Tuesday} \qquad\mathrm{(C)}\ \text{Wednesday} \qquad\mathrm{(D)}\ \text{Thursday} \qquad\mathrm{(E)}\ \text{Friday}$

 \textbf{Answer: }D



\textbf{Solution 1}

If there are five Mondays, there are only three possibilities for their dates: $(1,8,15,22,29)$, $(2,9,16,23,30)$, and $(3,10,17,24,31)$. 

 In the first case August starts on a Thursday, and there are five Thursdays, Fridays, and Saturdays in August. 

 In the second case August starts on a Wednesday, and there are five Wednesdays, Thursdays, and Fridays in August.

 In the third case August starts on a Tuesday, and there are five Tuesdays, Wednesdays, and Thursdays in August.

 The only day of the week that is guaranteed to appear five times is therefore $\boxed{\textrm{(D)}\ \text{Thursday}}$.

 



\textbf{Solution 2 (visualization)}

A diagram can be constructed to visualize the possibilities for the dates. First, a 4x7 rectangle representing 28 days can be constructed with three extra cells on the bottom to represent the three extra days. 

 Then, the rest of the steps in Solution 1 can be followed.

 For example, if the first of the bottom three cells is a Monday, it is clear that the following two days are Tuesday and Wednesday. This means that the new bottom three cells (as both months have 31 days) of August will be Thursday, Friday, and Saturday. This same process can be repeated so that the second cell is represented as Monday and so forth; then, the only day of the week that appears when all three cases are considered is $\boxed{\textrm{(D)}\ \text{Thursday}}$.

 


	\item (\textit{AMC 12 12B 2002}) If $a,b,c,d$ are positive real numbers such that $a,b,c,d$ form an increasing arithmetic sequence and $a,b,d$ form a geometric sequence, then $\frac ad$ is


 $\mathrm{(A)}\ \frac 1{12} \qquad\mathrm{(B)}\ \frac 16 \qquad\mathrm{(C)}\ \frac 14 \qquad\mathrm{(D)}\ \frac 13 \qquad\mathrm{(E)}\ \frac 12$

 \textbf{Answer: }C



\textbf{Solution 1}

We can let $a=1$, $b=2$, $c=3$, and $d=4$. $\frac{a}{d}=\boxed{\frac{1}{4}}  \Longrightarrow \mathrm{(C)}$




\textbf{Solution 2}

As $a, b, d$ is a geometric sequence, let $b=ka$ and $d=k^2a$ for some $k>0$. 

 Now, $a, b, c, d$ is an arithmetic sequence. Its difference is $b-a=(k-1)a$. Thus $d=a + 3(k-1)a = (3k-2)a$.

 Comparing the two expressions for $d$ we get $k^2=3k-2$. The positive solution is $k=2$, and $\frac{a}{d}=\frac{a}{k^2a}=\frac{1}{k^2}=\boxed{\frac{1}{4}} \Longrightarrow \mathrm{(C)}$.

 



\textbf{Solution 3}

Letting $n$ be the common difference of the arithmetic progression, we have$b = a + n$, $c = a + 2n$, $d = a + 3n$. We are given that $b / a$ = $d / b$, or
 \[\frac{a + n}{a} = \frac{a + 3n}{a + n}.\]Cross-multiplying, we get 
 \[a^2 + 2an + n^2 = a^2 + 3an\]
 \[n^2 = an\]
 \[n = a\]So $\frac{a}{d} = \frac{a}{a + 3n} = \frac{a}{4a} = \boxed{\frac{1}{4}}  \Longrightarrow \mathrm{(C)}$.

 


	\item (\textit{AMC 12 12B 2002}) How many different integers can be expressed as the sum of three distinct members of the set $\{1,4,7,10,13,16,19\}$?


 $\mathrm{(A)}\ 13 \qquad\mathrm{(B)}\ 16 \qquad\mathrm{(C)}\ 24 \qquad\mathrm{(D)}\ 30 \qquad\mathrm{(E)}\ 35$

 \textbf{Answer: }A



\textbf{Solution 1}

Subtracting 10 from each number in the set, and dividing the results by 3, we obtain the set $\{-3, -2, -1, 0, 1, 2, 3\}$.  It is easy to see that we can get any integer between $-6$ and $6$ inclusive as the sum of three elements from this set, for the total of $\boxed{\mathrm{(A) } 13}$ integers.

 



\textbf{Solution 2}

The set is an arithmetic sequence of numbers each $1$ more than a multiple of $3$. Thus the sum of any three numbers will be a multiple of $3$. All the multiples of $3$ from $1+4+7=12$ to $13+16+19=48$ are possible, totaling to $\boxed{\mathrm{(A) } 13}$ integers.

 


	\item (\textit{AMC 12 12B 2002}) The positive integers $A, B, A-B,$ and $A+B$ are all prime numbers. The sum of these four primes is


 $\mathrm{(A)}\ \mathrm{even} \qquad\mathrm{(B)}\ \mathrm{divisible\ by\ }3 \qquad\mathrm{(C)}\ \mathrm{divisible\ by\ }5 \qquad\mathrm{(D)}\ \mathrm{divisible\ by\ }7 \qquad\mathrm{(E)}\ \mathrm{prime}$

 \textbf{Answer: }E



\textbf{Solution 1}

Since $A-B$ and $A+B$ must have the same \href{https://artofproblemsolving.com/wiki/index.php?title=Parity}{parity}, and since there is only one even prime number, it follows that $A-B$ and $A+B$ are both odd. Thus one of $A, B$ is odd and the other even. Since $A+B > A > A-B > 2$, it follows that $A$ (as a prime greater than $2$) is odd. Thus $B = 2$, and $A-2, A, A+2$ are consecutive odd primes. At least one of $A-2, A, A+2$ is divisible by $3$, from which it follows that $A-2 = 3$ and $A = 5$. The sum of these numbers is thus $17$, a prime, so the answer is $\boxed{\mathrm{(E)}\ \text{prime}}$.

 



\textbf{Solution 2}

In order for both $A - B$ and $A + B$ to be prime, one of $A, B$ must be 2, or else both $A - B$, $A + B$ would be even numbers. 

 If $A = 2$, then $A < B$ and $A - B < 0$, which is not possible. Thus $B = 2$.

 Since $A$ is prime and $A > A - B > 2$, we can infer that $A > 3$ and thus $A$ can be expressed as $6n \pm 1$ for some natural number $n$. 

 However in either case, one of $A - B$ and $A + B$ can be expressed as $6n \pm 3 = 3(2n \pm 1)$ which is a multiple of 3. Therefore the only possibility that works is when $A - B = 3$ and 
 \[A + B + (A - B) + (A + B) = 5 + 2 + 3 + 7 = 17\]

 Which is a prime number. $\boxed{(E)}$




\textbf{Solution 3 (intuitive)}

Trying out some primes for $A$ and $B$ such that $A-B$ and $A+B$ are prime, $A=5$ and $B=2$ can be found almost immediately. Summing the four primes, the result is $5+2+3+7=17$, which is $\boxed{\mathrm{(E)}\ \text{prime}}$.

 


	\item (\textit{AMC 12 12B 2002}) For how many integers $n$ is $\dfrac n{20-n}$ the square of an integer?


 $\mathrm{(A)}\ 1 \qquad\mathrm{(B)}\ 2 \qquad\mathrm{(C)}\ 3 \qquad\mathrm{(D)}\ 4 \qquad\mathrm{(E)}\ 10$

 \textbf{Answer: }D



\textbf{Solution 1}

Let $x^2 = \frac{n}{20-n}$, with $x \ge 0$ (note that the solutions $x < 0$ do not give any additional solutions for $n$). Then rewriting, $n = \frac{20x^2}{x^2 + 1}$. Since $\text{gcd}(x^2, x^2 + 1) = 1$, it follows that $x^2 + 1$ divides $20$. Listing the factors of $20$, we find that $x = 0, 1, 2 , 3$ are the only $\boxed{\mathrm{(D)}\ 4}$ solutions (respectively yielding $n = 0, 10, 16, 18$).

 



\textbf{Solution 2}

For $n<0$ and $n>20$ the fraction is negative, for $n=20$ it is not defined, and for $n\in\{1,\dots,9\}$ it is between 0 and 1.

 Thus we only need to examine $n=0$ and $n\in\{10,\dots,19\}$.

 For $n=0$ and $n=10$ we obviously get the squares $0$ and $1$ respectively.

 For prime $n$ the fraction will not be an integer, as the denominator will not contain the prime in the numerator.

 This leaves $n\in\{12,14,15,16,18\}$, and a quick substitution shows that out of these only $n=16$ and $n=18$ yield a square. Therefore, there are only $\boxed{\mathrm{(D)}\ 4}$ solutions (respectively yielding $n = 0, 10, 16, 18$).

 



\textbf{Solution 3}

If $\frac{n}{20-n} = k^2 \ge 0$, then $n \ge 0$ and $20-n > 0$, otherwise $\frac{n}{20-n}$ will be negative. Thus $0 \le n \le 19$ and 
 \[0 = \frac{0}{20-(0)} \le \frac{n}{20-n} \le \frac{19}{20-(19)} = 19\] Checking all $k$ for which $0 \le k^2 \le 19$, we have $0$, $1$, $2$, $3$ as the possibilities. $\boxed{(D)}$




\textbf{Solution 4}

For all integers x, $x^2$ is always a positive integer. So solve for $\frac{n}{20-n} = 0$, getting $n=0$ and $\frac{n}{20-n} = 1$, getting $n =10$. For all values n less than 0 and greater than 20, the value $\frac{n}{20-n}$ is negative, so now try values of n between 10 and 20. Quick substitution finds $0$, $10$, $16$, and $18$ which yields $x=0$, $x=1$, $x=2$, and $x=3$ respectively. 4 values, or $\boxed{(D)}$




\textbf{Solution 5}

Simon's Favourite Factoring Trick.

 Since $\frac{n}{20-n}$ is an integer $k$, we multiply both sides by $20-n$. This gives us $n=20k^2-nk^2$. We subtract $20k^2$ on both sides, then add $nk^2$ on both sides as a prerequisite for using Simon's Favorite Factoring Trick. We have $(k^2+1)(n-20)=20$. We then consider the different factors of $20$ that $k^2+1$ can be. It could be $1,2,4,5,10$, and $20$. After checking case by case, we then are able to identify that there are $4$ such $k$ values that also yield an integer $n$ value, meaning that there are $4$ values, so the correct answer is  $\boxed{(D)}$~CharmaineMa0729201

 



\textbf{Solution 6}

Similar to previous solutions, we can bound $n$ between $0$ and $20$. Then, we let $\frac{n}{20-n} = k^2$ for some integer $k$. Multiplying out the fraction and factoring, we get $n(k^2 + 1) = k^2 \cdot 20$. As before, $\text{gcd}(k^2 + 1, k^2) = 1$, so we have that $k^2 | n$. This means that $0 \leq k^2 \leq 20$, as $k^2$ now has to be at most $n$, so trying $k^2 = 0, 1, 4, 9, 16$, we get four possible $n$, namely $0$, $10$, $16$, and $18$. Thus, $\boxed{(D)}$.

 


	\item (\textit{AMC 12 12B 2002}) The sum of $18$ consecutive positive integers is a perfect square. The smallest possible value of this sum is


 $\mathrm{(A)}\ 169 \qquad\mathrm{(B)}\ 225 \qquad\mathrm{(C)}\ 289 \qquad\mathrm{(D)}\ 361 \qquad\mathrm{(E)}\ 441$

 \textbf{Answer: }B



\textbf{Solution 1}

Let $a, a+1, \ldots, a + 17$ be the consecutive positive integers. Their sum, $18a + \frac{17(18)}{2} = 9(2a+17)$, is a perfect square. Since $9$ is a perfect square, it follows that $2a + 17$ is a perfect square. The smallest possible such perfect square is $25$ when $a = 4$, and the sum is $225 \Rightarrow \mathrm{(B)}$.

 



\textbf{Solution 2}

Notice that all five choices given are perfect squares.

 Let $a$ be the smallest number, we have 
 \[a+(a+1)+(a+2)+...+(a+17)=18a+\sum_{k=1}^{17}k=18a+153\]

 Subtract $153$ from each of the choices and then check its divisibility by $18$, we have $225$ as the smallest possible sum. $\mathrm {(B)}$




\textbf{Solution 1.1}

the normal sequence can be described as $N^2+N$ divided by 2.

 Since have 18 terms adding 18n will increase the consective sequence startering number by 1 

 $\frac{(N^2+N)}{2} +18n$

 now subsitute 18 as N

 we get $\frac{(18^2+18)}{2} = 171$

 
put $I^2$ which is integer square and plug in all our results $171 + 18n = I^2$
$18n = I^2-171 I^2-171 = 0 \pmod{18}$

 subsitute the answer choices starting with B because 169 is less than 171 and results in a neagtive number

 $225-177 = \pmod{18}$$54 = \pmod{18}$

 54 is divisble by 18 and is therefore the smallest number possible.

 


	\item (\textit{AMC 12 12B 2002}) Four distinct circles are drawn in a plane. What is the maximum number of points where at least two of the circles intersect?


 $\mathrm{(A)}\ 8 \qquad\mathrm{(B)}\ 9 \qquad\mathrm{(C)}\ 10 \qquad\mathrm{(D)}\ 12 \qquad\mathrm{(E)}\ 16$

 \textbf{Answer: }D



\textbf{Solution 1}

For any given pair of circles, they can intersect at most $2$ times. Since there are ${4\choose 2} = 6$ pairs of circles, the maximum number of possible intersections is $6 \cdot 2 = 12$. We can construct such a situation as below, so the answer is $\boxed{\mathrm{(D)}\ 12}$.

 
                    \begin{figure}[H]
                        \centering
                        \includegraphics[width=0.5\textwidth]{Solution/AMC_12_12B_2002_14_Solution_Figure_1.png}
                    \end{figure}
                    




\textbf{Solution 2}

Because a pair or circles can intersect at most $2$ times, the first circle can intersect the second at $2$ points, the third can intersect the first two at $4$ points, and the fourth can intersect the first three at $6$ points. This means that our answer is $2+4+6=\boxed{\mathrm{(D)}\ 12}.$




\textbf{Solution 3}

Pick a circle any circle- $4$ ways. Then, pick any other circle- $3$ ways. For each of these circles, there will be $2$ intersections for a total of $4*3*2$ = $24$ intersections. However, we have counted each intersection twice, so we divide for overcounting. Therefore, we reach a total of $\frac{24}{2}=\boxed{12}$, which corresponds to $\text{(D)}$.

 


	\item (\textit{AMC 12 12B 2002}) How many four-digit numbers $N$ have the property that the three-digit number obtained by removing the leftmost digit is one ninth of $N$?


 $\mathrm{(A)}\ 4 \qquad\mathrm{(B)}\ 5 \qquad\mathrm{(C)}\ 6 \qquad\mathrm{(D)}\ 7 \qquad\mathrm{(E)}\ 8$

 \textbf{Answer: }D



\textbf{Solution}

Let $N = \overline{abcd} = 1000a + \overline{bcd}$, such that $\frac{N}{9} = \overline{bcd}$. Then $1000a + \overline{bcd} = 9\overline{bcd} \Longrightarrow 125a = \overline{bcd}$. Since $100 \le \overline{bcd} < 1000$, from $a = 1, \ldots, 7$ we have $7$ three-digit solutions, and the answer is  $\mathrm{(D)}$. 

 


 



\textbf{Solution 1}

Since N is a four digit number, assume WLOG that $N = 1000a + 100b + 10c + d$, where a is the thousands digit, b is the hundreds digit, c is the tens digit, and d is the ones digit.Then, $\frac{1}{9}N = 100b + 10c + d$, so $N = 900b + 90c + 9d$Set these equal to each other:
 \[1000a + 100b + 10c + d = 900b + 90c + 9d\]
 \[1000a = 800b + 80c + 8d\]
 \[1000a = 8(100b + 10c + d)\]Notice that $100b + 10c + d = N - 1000a$, thus:
 \[1000a = 8(N - 1000a)\]
 \[1000a = 8N - 8000a\]
 \[9000a = 8N\]
 \[N = 1125a\]

 Go back to our first equation, in which we set $N = 1000a + 100b + 10c + d$,Then: 
 \[1125a = 1000a + 100b + 10c + d\]
 \[125a = 100b + 10c + d\]The upper limit for the right hand side (RHS) is $999$ (when $b = 9$, $c = 9$, and $d = 9$).It's easy to prove that for an $a$ there is only one combination of $b, c,$ and $d$ that can make the equation equal. Just think about the RHS as a three digit number $bcd$. There's one and only one way to create every three digit number with a certain combination of digits.Thus, we test for how many as are in the domain set by the RHS. Since $125\cdot7 = 875$ which is the largest $a$ value, then $a$ can be $1$ through $7$, giving us the answer of $\boxed {D) 7}$

 IronicNinja~

 



\textbf{Solution 2}

Let Let $a$ denote the leftmost digit of $N$ and let $x$ denote the three-digit number obtained by removing $a$. Then $N=1000a+x=9x$ and it follows that $1000a=8x$. Dividing both sides by 8 yields $125a=x$. All the values of $a$ in the range 1 to 7 result in three-digit numbers, hence there are $\boxed{D)7}$ values for $N$.

 $\sim$Glowworm

 



\textbf{Solution 3 (cases)}

Let $a, b, c, d$ denote the number $\overline{abcd}$. If $a = 1$, $\overline{abcd} - 1000 = \frac{\overline{abcd}}{9}$. Clearly, $\overline{abcd} = 1125$. Doing the same for 2 yields $2250$, and 3 gives $3375$. Clearly there is a pattern of adding 125, which works up until 7, after which we get $9000$ which clearly doesn't work. Therefore, there are $\boxed {D) 7}$ valid solutions, those being:$1125, 2250, 3375, 4500, 5625, 6750, 7875$.

 


	\item (\textit{AMC 12 12B 2002}) Juan rolls a fair regular octahedral die marked with the numbers $1$ through $8$. Then Amal rolls a fair six-sided die. What is the probability that the product of the two rolls is a multiple of 3?


 $\mathrm{(A)}\ \frac1{12} \qquad\mathrm{(B)}\ \frac 13 \qquad\mathrm{(C)}\ \frac 12 \qquad\mathrm{(D)}\  \frac 7{12} \qquad\mathrm{(E)}\ \frac 23$

 \textbf{Answer: }C



\textbf{Solution 1}

On both dice, only the faces with the numbers $3,6$ are divisible by $3$. Let $P(a) = \frac{2}{8} = \frac{1}{4}$ be the probability that Juan rolls a $3$ or a $6$, and $P(b) = \frac{2}{6} = \frac 13$ that Amal does. By the \href{https://artofproblemsolving.com/wiki/index.php?title=Principle\_of\_Inclusion-Exclusion}{Principle of Inclusion-Exclusion},

 
 \[P(a \cup b) = P(a) + P(b) - P(a \cap b) = \frac{1}{4} + \frac{1}{3} - \frac{1}{4} \cdot \frac{1}{3} = \frac{1}{2} \Rightarrow \mathrm{(C)}\]

 Alternatively, the probability that Juan rolls a multiple of $3$ is $\frac{1}{4}$, and the probability that Juan does not roll a multiple of $3$ but Amal does is $\left(1 - \frac{1}{4}\right) \cdot \frac{1}{3} = \frac{1}{4}$. Thus the total probability is $\frac 14 + \frac 14 = \frac 12$.

 



\textbf{Solution 2}

The probability that neither Juan nor Amal rolls a multiple of $3$ is $\frac{6}{8} \cdot \frac{4}{6} = \frac{1}{2}$; using \href{https://artofproblemsolving.com/wiki/index.php?title=Complementary\_counting}{complementary counting}, the probability that at least one does is $1 - \frac 12 = \frac 12 \Rightarrow \mathrm{(C)}$.

 



\textbf{Solution 3 (Alcumus)}

The product will be a multiple of 3 if and only if at least one of the two rolls is a 3 or a 6. The probability that Juan rolls 3 or 6 is $2/8 = 1/4$. The probability that Juan does not roll 3 or 6, but Amal does is $(3/4) (1/3) = 1/4$. Thus, the probability that the product of the rolls is a multiple of 3 is 
 \[\frac{1}{4} + \frac{1}{4} = \boxed{\frac{1}{2}}.\]~aopsav (Credit to AoPS Alcumus)

 


	\item (\textit{AMC 12 12B 2002}) Andy’s lawn has twice as much area as Beth’s lawn and three times as much area as Carlos’ lawn. Carlos’ lawn mower cuts half as fast as Beth’s mower and one third as fast as Andy’s mower. If they all start to mow their lawns at the same time, who will finish first?


 $\mathrm{(A)}\ \text{Andy} \qquad\mathrm{(B)}\ \text{Beth} \qquad\mathrm{(C)}\ \text{Carlos} \qquad\mathrm{(D)}\ \text{Andy\ and \ Carlos\ tie\ for\ first.} \qquad\mathrm{(E)}\ \text{All\ three\ tie.}$

 \textbf{Answer: }B



\textbf{Solution 1}

We say Andy's lawn has an area of $x$. Beth's lawn thus has an area of $\frac{x}{2}$, and Carlos's lawn has an area of $\frac{x}{3}$. 

 We say Andy's lawn mower cuts at a speed of $y$. Carlos's cuts at a speed of $\frac{y}{3}$, and Beth's cuts at a speed $\frac{2y}{3}$.

 Each person's lawn is cut at a time of $\frac{\text{area}}{\text{rate}}$, so Andy's is cut in $\frac{x}{y}$ time, as is Carlos's. Beth's is cut in $\frac{3}{4}\times\frac{x}{y}$, so the first one to finish is $\boxed{\mathrm{(B)}\ \text{Beth}}$.

 



\textbf{Solution 2}

WLOG, we can set values of their lawns' areas and their owners' speeds. Let the area of Andy's lawn be $6$ units, Beth's lawn be $3$ units, and Carlos's lawn be $2$ units. Let Carlos's mowing area per hour (honestly the time you set won't matter) be $1$, let Beth's mowing area per hour be $2$, and let Andy's mowing area per hour be $3$. Now, we can easily calculate their time  by dividing their lawns' areas by their respective owners' speeds. Andy's time is $\frac{6}{3} = 2$ hours, Beth's time is $\frac{3}{2} = 1.5$ hours, and Carlos's time is $\frac{6}{3} = 2$ hours. Our answer is clearly $\boxed{\mathrm{(B)}\ \text{Beth}}$.~Solution by virjoy2001

 



\textbf{Solution 3(convert into a racing problem)}

We first set the area for Carlo's lawn to $x$. That leads to Beth's lawn's area to be $3x$ and Andy's lawn $6x$. After that, we set the speed of Carlo's lawn mower to $y$. We then get Beth's speed in terms of $y$, which is $2y$. Andy's lawn mower's speed is $3y$. WLOG, we set the width of their lawns to $0$. We then can imagine that it's a racetrack, with the distance being the area of their lawns, and they are competing with different speeds and distances. We then proceed by applying the speed formula, distance = speed*time. Using WLOG again, we can set Carlo's speed to $10$ mph and we can assume Carlo's lawn, leaving Beth with $20$ mph/$3$miles and Andy with $30$ mph/$6$ miles. After plugging in the variables, we clearly see that $\boxed{\mathrm{(B)}\ \text{Beth}}$ will win.~Calvin67

 


	\item (\textit{AMC 12 12B 2002}) A point $P$ is randomly selected from the rectangular region with vertices $(0,0),(2,0),(2,1),(0,1)$. What is the probability that $P$ is closer to the origin than it is to the point $(3,1)$?


 

$\mathrm{(A)}\ \frac 12 \qquad\mathrm{(B)}\ \frac 23 \qquad\mathrm{(C)}\ \frac 34 \qquad\mathrm{(D)}\ \frac 45 \qquad\mathrm{(E)}\ 1$

 \textbf{Answer: }C



\textbf{Solution 1}

Assume that the point $P$ is randomly chosen within the rectangle with vertices $(0,0)$, $(3,0)$, $(3,1)$, $(0,1)$. In this case, the region for $P$ to be closer to the origin than to point $(3,1)$ occupies exactly $\frac{1}{2}$ of the area of the rectangle, or $1.5$ square units.

 
If $P$ is chosen within the square with vertices $(2,0)$, $(3,0)$, $(3,1)$, $(2,1)$ which has area $1$ square unit, it is for sure closer to $(3,1)$.

 
Now if $P$ can only be chosen within the rectangle with vertices $(0,0)$, $(2,0)$, $(2,1)$, $(0,1)$, then the square region is removed and the area for $P$ to be closer to $(3,1)$ is then decreased by $1$ square unit, left with only $0.5$ square unit.

 
Thus the probability that $P$ is closer to $(3.1)$ is $\frac{0.5}{2}=\frac{1}{4}$ and that of $P$ is closer to the origin is $1-\frac{1}{4}=\frac{3}{4}$. $\mathrm{(C)}$




\textbf{Solution 2 (slightly more calculation but easy:)}

First, join points $(0,0)$ and $(3,1)$. This line $l_1$ has equation $y = \frac{1}{3}x$. 

 
Next, consider the perpendicular bisector $l_2$ of line $l_1$, any point on the perpendicular bisector is equidistance from points $(0,0)$ and $(3,1)$. 

 
If the point $P$ is chosen on the left of $l_2$ , the point is closer to the origin. $l_2$ will cut the rectangle region twice, dividing the region into two smaller trapezoids. The trapezoid on the left side is the area we want. 

 
So, first, find the equation of $l_2$, using the midpoint of $(0,0)$ and $(3,1)$, which is $(1.5,0.5)$, the equation $y = -3x + 5$ can be easily derived. Substituting $y = 0$ and $y = 1$ to solve for $x$. For the former, $x = \frac{5}{3}$, for the latter, $x = \frac{4}{3}$.

 
Thus, the area of the trapezoid we want is $\frac {(\frac{5}{3} + \frac{4}{3})\cdot1}{2} = \frac{3}{2}$

 Therefore, the probability the question asks equals to $\frac{\frac{3}{2}}{2} =\boxed{\textbf{(C) }\frac{3}{4}}$



	\item (\textit{AMC 12 12B 2002}) If $a,b,$ and $c$ are positive real numbers such that $a(b+c) = 152, b(c+a) = 162,$ and $c(a+b) = 170$, then $abc$ is


 $\mathrm{(A)}\ 672 \qquad\mathrm{(B)}\ 688 \qquad\mathrm{(C)}\ 704 \qquad\mathrm{(D)}\ 720 \qquad\mathrm{(E)}\ 750$

 \textbf{Answer: }D



\textbf{Solution 1}

Adding up the three equations gives $2(ab + bc + ca) = 152 + 162 + 170 = 484 \Longrightarrow ab + bc + ca = 242$. Subtracting each of the above equations from this yields, respectively, $bc = 90, ca = 80, ab = 72$. Taking their product, $ab \cdot bc \cdot ca = a^2b^2c^2 = 90 \cdot 80 \cdot 72 = 720^2 \Longrightarrow abc = \boxed{720} \Rightarrow \mathrm{(D)}$.

 



\textbf{Solution 2}

Expand the system and subtract equation 3 from equation 2 to get $ab-ac=-8$, and add it to the first equation to get $2ab=144$ $\Longrightarrow$ $ab=72.$ Since we want $abc,$ it's trivial that 72 is a factor of this. We can quickly see that none of the answer choices have the factor $72$ except choice $\boxed{\mathrm{(D)}}.$



	\item (\textit{AMC 12 12B 2002}) Let $\triangle XOY$ be a right-angled triangle with $\angle XOY = 90^{\circ}$. Let $M$ and $N$ be the midpoints of legs $OX$ and $OY$, respectively. Given that $XN = 19$ and $YM = 22$, find $XY$.


 $\mathrm{(A)}\ 24 \qquad\mathrm{(B)}\ 26 \qquad\mathrm{(C)}\ 28 \qquad\mathrm{(D)}\ 30 \qquad\mathrm{(E)}\ 32$

 \textbf{Answer: }B



\textbf{Solution 1}


                    \begin{figure}[H]
                        \centering
                        \includegraphics[width=0.5\textwidth]{Solution/AMC_12_12B_2002_20_Solution_Figure_1.png}
                    \end{figure}
                    

 Let $OM = x$, $ON = y$. By the \href{https://artofproblemsolving.com/wiki/index.php?title=Pythagorean\_Theorem}{Pythagorean Theorem} on $\triangle XON, MOY$ respectively,
 \begin{align*} (2x)^2 + y^2 &= 19^2\\ x^2 + (2y)^2 &= 22^2\end{align*}


 Summing these gives $5x^2 + 5y^2 = 845 \Longrightarrow x^2 + y^2 = 169$. 

 By the Pythagorean Theorem again, we have 

 
 \[(2x)^2 + (2y)^2 = XY^2 \Longrightarrow XY = \sqrt{4(x^2 + y^2)} = \sqrt{4(169)} = \sqrt{676} = \boxed{\mathrm{(B)}\ 26}.\]

 Alternatively, we could note that since we found $x^2 + y^2 = 169$, segment $MN=13$. Right triangles $\triangle MON$ and $\triangle XOY$ are similar by Leg-Leg with a ratio of $\frac{1}{2}$, so $XY=2(MN)=\boxed{\mathrm{(B)}\ 26}$.

 


 



\textbf{Solution 2}

Let $XO=x$ and $YO=y.$ Then, $XY=\sqrt{x^2+y^2}.$ 

 Since $XN=19$ and $YM=22,$ 
 \[XN^2=19^2=x^2+(\dfrac{y}{2})^2)=\dfrac{x^2}{4}+y^2\]
 \[YM^2=22^2=(\dfrac{x}{2})^2+y^2=\dfrac{x^2}{4}+y^2.\]

 Adding these up: 
 \[19^2+22^2=\dfrac{4x^2+y^2}{4}+\dfrac{x^2+4y^2}{4}\]
 \[845=\dfrac{5x^2+5y^2}{4}\]
 \[3380=5x^2+5y^2\]
 \[676=x^2+y^2.\] 

 Then, we substitute: $XY=\sqrt{x^2+y^2}=\sqrt{676}=\boxed{26}.$




\textbf{Solution 3 (Solution 1 but shorter)}

Refer to the diagram in solution 1. $4x^2+y^2=361$ and $4y^2+x^2=484$, so add them: $5x^2+5y^2=845$ and divide by 5: $x^2+y^2=169$ so $\dfrac{XY}{2}=\sqrt{169}=13$ and so $XY=26$, or answer $\boxed{\text{(B)}}$.

 



\textbf{Solution 4}

Use the diagram in solution 1. Get $4x^2+y^2=361$ and $4y^2+x^2=484$, and multiply the second equation by 4 to get $4x^2+16y^2=1936$ and then subtract the first from the second. Get $15y^2=1575$ and $y^2=105$. Repeat for the other variable to get $15x^2=960$ and $x^2=64$. Now XY is equal to the square root of four times these quantities, so $(105+64) \cdot 4=676$, and $XY=\boxed{\text{(B)}\ 26}$.

 



\textbf{Solution 5 (Quick and good for in-contest use)}

Let $XM = MO = y$, and $ON = NY  = x$. Now, we can write some equations. We have that 
 \[(2y)^2+x^2 = 19^2 \Longrightarrow 4y^2 + x^2 = 361,\] and 
 \[(2x)^2+y^2 = 22^2 \Longrightarrow 4x^2+y^2 = 484.\] Adding the $2$ equations up we get 
 \[5x^2+5y^2 = 845.\] The hypotenuse of the triangle is equal to 
 \[\sqrt{4x^2+4y^2},\] so that is what we want. In order to get 
 \[4x^2+4y^2,\] we need to subtract $x^2+y^2$ from $5x^2+5y^2$, and since $x^2+y^2 = \dfrac{1}{5}\left(5x^2+5y^2\right)$, we have that 
 \[4x^2+4y^2 = \dfrac{4}{5} \cdot 845 = 676.\] Therefore, the hypotenuse is $\sqrt{676} = \boxed{\text{(B)}\ 26}$.

 -jb2015007

 


	\item (\textit{AMC 12 12B 2002}) For all positive integers $n$ less than $2002$, let 


 
 \begin{eqnarray*} a_n =\left\{ \begin{array}{lr} 11, & \text{if\ }n\ \text{is\ divisible\ by\ }13\ \text{and\ }14;\\ 13, & \text{if\ }n\ \text{is\ divisible\ by\ }14\ \text{and\ }11;\\ 14, & \text{if\ }n\ \text{is\ divisible\ by\ }11\ \text{and\ }13;\\ 0, & \text{otherwise}. \end{array} \right. \end{eqnarray*}


 Calculate $\sum_{n=1}^{2001} a_n$.


 $\mathrm{(A)}\ 448 \qquad\mathrm{(B)}\ 486 \qquad\mathrm{(C)}\ 1560 \qquad\mathrm{(D)}\ 2001 \qquad\mathrm{(E)}\ 2002$

 \textbf{Answer: }A



\textbf{Solution 1}

Since $2002 = 11 \cdot 13 \cdot 14$, it follows that 
 \begin{eqnarray*} a_n =\left\{ \begin{array}{lr} 11, & \text{if\ }n=13 \cdot 14 \cdot k, \quad k = 1,2,\cdots 10;\\ 13, & \text{if\ }n=14 \cdot 11 \cdot k, \quad k = 1,2,\cdots 12;\\ 14, & \text{if\ }n=11 \cdot 13 \cdot k, \quad k = 1,2,\cdots 13;\\ \end{array} \right. \end{eqnarray*}

 Thus $\sum_{n=1}^{2001} a_n = 11 \cdot 10 + 13 \cdot 12 + 14 \cdot 13 = 448 \Rightarrow \mathrm{(A)}$.$\begin{array}{lr} 11, & \text{if\ }n=13 \cdot 14 \cdot k, \quad k = 1,2,\cdots 10;\\ 13, & \text{if\ }n=14 \cdot 11 \cdot k, \quad k = 1,2,\cdots 12;\\ 14, & \text{if\ }n=11 \cdot 13 \cdot k, \quad k = 1,2,\cdots 13;\\ \end{array}$.

 



\textbf{Solution 2}

Find the LCMs of the groups of the numbers.

 Notice that the groups are relatively prime.

 So $a_n=$:

 11 if $n$ is a multiple of 182.

 13 if $n$ is a multiple of 154.

 14 if $n$ is a multiple of 143.

 When do we see ambiguities (for example: $n$ is a multiple of 11, 13, and 14)? This is only done when $n$ is a multiple of $\operatorname{lcm}(11,13,14)=2002$. However, since $n<2002$, this can never happen.

 So we have 10 multiples of 182 we have to count (1 to 10 $*182$), and similarly, 12 multiples of 154, and 13 multiples of 143. The sum is $10*11+12*13+13*14=448$. Select $\boxed{A}$.

 


	\item (\textit{AMC 12 12B 2002}) For all integers $n$ greater than $1$, define $a_n = \frac{1}{\log_n 2002}$. Let $b = a_2 + a_3 + a_4 + a_5$ and $c = a_{10} + a_{11} + a_{12} + a_{13} + a_{14}$. Then $b- c$ equals


 $\mathrm{(A)}\ -2 \qquad\mathrm{(B)}\ -1  \qquad\mathrm{(C)}\ \frac{1}{2002} \qquad\mathrm{(D)}\ \frac{1}{1001} \qquad\mathrm{(E)}\ \frac 12$

 \textbf{Answer: }B



\textbf{Solution 1}

By the \href{https://artofproblemsolving.com/wiki/index.php?title=Logarithms\#Logarithmic\_Properties}{change of base} formula, $a_n = \frac{1}{\frac{\log 2002}{\log n}} = \left(\frac{1}{\log 2002}\right) \log n$. Thus 
 \begin{align*}b- c &= \left(\frac{1}{\log 2002}\right)(\log 2 + \log 3 + \log 4 + \log 5 - \log 10 - \log 11 - \log 12 - \log 13 - \log 14)\\ &= \left(\frac{1}{\log 2002}\right)\left(\log \frac{2 \cdot 3 \cdot 4 \cdot 5}{10 \cdot 11 \cdot 12 \cdot 13 \cdot 14}\right)\\  &= \left(\frac{1}{\log 2002}\right) \log 2002^{-1} = -\left(\frac{\log 2002}{\log 2002}\right) = -1 \Rightarrow \mathrm{(B)}\end{align*}





\textbf{Solution 2 (basically solution 3 but with steps and a quicker conclusion)}

Note that \( \frac{1}{\log_a(b)} = \log_b(a) \). Therefore, \( \frac{1}{\log_n(2002)} = \log_{2002}(n) \). For the sake of space and time, set \( a = 2002 \). Then \( b \) can be represented as
 \[\log_a(2) + \log_a(3) + \log_a(4) + \log_a(5)\]
 \[\text{which is simply}\]
 \[\log_a(2 \cdot 3 \cdot 4 \cdot 5)\]
 \[\text{and } c \text{ can be represented as}\]
 \[\log_a(10) + \log_a(11) + \log_a(12) + \log_a(13) + \log_a(14)\]
 \[\text{which is simply}\]
 \[\log_a(10 \cdot 11 \cdot 12 \cdot 13 \cdot 14)\]
 \[\text{so } b - c \text{ is just}\]
 \[\log_a(2 \cdot 3 \cdot 4 \cdot 5) - \log_a(10 \cdot 11 \cdot 12 \cdot 13 \cdot 14)\]
 \[\text{which is simply}\]
 \[\log_a\left(\frac{2 \cdot 3 \cdot 4 \cdot 5}{10 \cdot 11 \cdot 12 \cdot 13 \cdot 14}\right)\]
 \[\text{or}\]
 \[\log_a\left(\frac{1}{2002}\right)\]
 \[\text{but } a = 2002, \text{ so we have}\]
 \[\log_{2002}\left(\frac{1}{2002}\right)\]
 \[=\boxed{\textbf{(B)}-1}\]




\textbf{Solution 3}

Note that $\frac{1}{\log_a b}=\log_b a$. Thus $a_n=\log_{2002} n$. Also notice that if we have a log sum, we multiply, and if we have a log product, we divide. Using these properties, we get that the result is the following:

 
 \[\log_{2002}\left(\frac{2*3*4*5}{10*11*12*13*14}=\frac{1}{11*13*14}=\frac{1}{2002}\right)=\boxed{\textbf{(B)}-1}\]




\textbf{Solution 4}

Note that $a_2 = \frac{1}{\log_2 2002}$. $1$ is also equal to $\log_2 2$. So $a_2 = \frac{\log_2 2}{\log_2 2002}$. By the change of bases formula, $a_2 = \log_{2002} 2$. Following the same reasoning, $a_3 = \log_{2002} 3$, $a_4 = \log_{2002} 4$ and so on.

 
 \[b = \log_{2002} 2 + \log_{2002} 3 + .....+ \log_{2002} 5 = \log_{2002} 5! = \log_{2002} 120\]Now solving for $c$, we see that it equals $\log_{2002} (10\cdot 11 \cdot 12 \cdot 13 \cdot 14)$
 \[b-c = \log_{2002} 120 - \log_{2002} 240240 \rightarrow \log_{2002} \frac{1}{2002} = \boxed{-1}\]



	\item (\textit{AMC 12 12B 2002}) In $\triangle ABC$, we have $AB = 1$ and $AC = 2$. Side $\overline{BC}$ and the median from $A$ to $\overline{BC}$ have the same length. What is $BC$?


 $\mathrm{(A)}\ \frac{1+\sqrt{2}}{2} \qquad\mathrm{(B)}\ \frac{1+\sqrt{3}}2 \qquad\mathrm{(C)}\ \sqrt{2} \qquad\mathrm{(D)}\ \frac 32 \qquad\mathrm{(E)}\ \sqrt{3}$

 \textbf{Answer: }C



\textbf{Solution 1: Pythagorean Theorem}


                    \begin{figure}[H]
                        \centering
                        \includegraphics[width=0.5\textwidth]{Solution/AMC_12_12B_2002_23_Solution_Figure_1.png}
                    \end{figure}
                    

 Let $D$ be the foot of the altitude from $A$ to $\overline{BC}$ extended past $B$. Let $AD = x$ and $BD = y$.Using the Pythagorean Theorem, we obtain the equations

 
 \begin{align*}  x^2 + y^2 = 1 \hspace{0.5cm}(1)\\ x^2 + y^2 + 2ya + a^2 = 4a^2 \hspace{0.5cm}(2)\\ x^2 + y^2 + 4ya + 4a^2 = 4 \hspace{0.5cm}(3) \end{align*}


 Subtracting $(1)$ equation from $(2)$ and $(3)$, we get

 
 \begin{align*} 2ya + a^2 = 4a^2 - 1 \hspace{0.5cm}(4)\\ 4ya + 4a^2 = 3 \hspace{0.5cm}(5) \end{align*}


 Then, subtracting $2 \times (4)$ from $(5)$ and rearranging, we get $10a^2 = 5$, so $BC = 2a = \sqrt{2}\Rightarrow \boxed{\mathrm{(C)}}$




\textbf{Solution 2: Law of Cosines}


                    \begin{figure}[H]
                        \centering
                        \includegraphics[width=0.5\textwidth]{Solution/AMC_12_12B_2002_23_Solution_Figure_2.png}
                    \end{figure}
                    

 Let $D$ be the foot of the median from $A$ to $\overline{BC}$, and we let $AD = BC = 2a$. Then by the \href{https://artofproblemsolving.com/wiki/index.php?title=Law\_of\_Cosines}{Law of Cosines} on $\triangle ABD, \triangle ACD$, we have 
 \begin{align*} 1^2 &= a^2 + (2a)^2 - 2(a)(2a)\cos ADB \\ 2^2 &= a^2 + (2a)^2 - 2(a)(2a)\cos ADC  \end{align*}


 Since $\cos ADC = \cos (180 - ADB) = -\cos ADB$, we can add these two equations and get 

 
 \[5 = 10a^2\]

 Hence $a = \frac{1}{\sqrt{2}}$ and $BC = 2a = \sqrt{2} \Rightarrow \mathrm{(C)}$.

 



\textbf{Solution 3: Stewart's Theorem}

From \href{https://artofproblemsolving.com/wiki/index.php?title=Stewart\%27s\_Theorem}{Stewart's Theorem}, we have $(2)(1/2)a(2) + (1)(1/2)a(1) = (a)(a)(a) + (1/2)a(a)(1/2)a.$ Simplifying, we get $(5/4)a^3 = (5/2)a \implies (5/4)a^2 = 5/2 \implies a^2 = 2 \implies a = \boxed{\sqrt{2}}.$- awu2014

 



\textbf{Solution 4: Pappus's Median Theorem}

There is a theorem in geometry known as Pappus's Median Theorem. It states that if you have $\triangle{ABC}$, and you draw a median from point $A$ to side $BC$ (label this as $M$), then: $(AM)^2 = \dfrac{2(b^2) + 2(c^2) - (a^2)}{4}$. Note that $b$ is the length of side $\overline{AC}$, $c$ is the length of side $\overline{AB}$, and $a$ is length of side $\overline{BC}$. Let $MB = MC = x$. Then $AM = 2x$. Now, we can plug into the formula given above: $AM = 2x$, $b = 2$, $c = 1$, and $a = 2x$. After some simple algebra, we find $x = \dfrac{\sqrt{2}}{2}$. Then, $BC = \boxed{\sqrt{2}} \implies \boxed{C}$. 

 -Flames

 Note: Pappus's Median Theorem is just a special case of Stewart's Theorem, with $m = n$. ~Puck\_0aka Apollonius' Theorem - Orion 2010

 



\textbf{Solution 5: Similar Triangles}


                    \begin{figure}[H]
                        \centering
                        \includegraphics[width=0.5\textwidth]{Solution/AMC_12_12B_2002_23_Solution_Figure_2.png}
                    \end{figure}
                    

 First let us call the intersection between $\overline{BC}$ and the median $M$.Notice that $\overline{AC}:\overline{AB}=2:1$. Also, $\overline{MA}:\overline{MB}=2:1$, too. Since $\triangle{ABC}$ and $\triangle{MBA}$ share the same angle $\angle{ABC}$, we can say that the triangles have SSA similarity.

 Technically, SSA similarity is invalid since there are two possible triangles which can be drawn with the same SSA. Just by looking at the diagram, we can tell that we have the correct triangle which is truly similar. If tried the same similarity with triangle $\triangle{MCA}$, it would not work.

 Now that we have similar triangles, we can say that $\overline{AB}:\overline{BM}=\overline{CB}:\overline{BA}$, so $1:a=2a:1$. This gives $2a^2=1$ and $a = \frac{\sqrt2}{2}$. Since the median has length $2a$, the final answer is $2\cdot\frac{\sqrt2}{2}=\boxed{\sqrt2}\implies \boxed{C}$.

 -mathisfunsp

 


	\item (\textit{AMC 12 12B 2002}) A convex quadrilateral $ABCD$ with area $2002$ contains a point $P$ in its interior such that $PA = 24, PB = 32, PC = 28, PD = 45$. Find the perimeter of $ABCD$.


 $\mathrm{(A)}\ 4\sqrt{2002} \qquad\mathrm{(B)}\ 2\sqrt{8465} \qquad\mathrm{(C)}\ 2(48+$ $\sqrt{2002}) \qquad\mathrm{(D)}\ 2\sqrt{8633} \qquad\mathrm{(E)}\ 4(36 + \sqrt{113})$

 \textbf{Answer: }E



\textbf{Solution 1}

We have 
 \[[ABCD] = 2002 \le \frac 12 (AC \cdot BD)\](This is true for any convex quadrilateral: split the quadrilateral along $AC$ and then using the triangle area formula to evaluate $[ACB]$ and $[ACD]$), with equality only if $\overline{AC} \perp \overline{BD}$. By the \href{https://artofproblemsolving.com/wiki/index.php?title=Triangle\_inequality}{triangle inequality},  

 
 \begin{align*}AC &\le PA + PC = 52\\ BD &\le PB + PD = 77\end{align*}


 with equality if $P$ lies on $\overline{AC}$ and $\overline{BD}$ respectively. Thus

 
 \[2002 \le \frac{1}{2} AC \cdot BD \le \frac 12 \cdot 52 \cdot 77 = 2002\]

 Since we have the equality case, $\overline{AC} \perp \overline{BD}$ at point $P$, as shown below. 

 
                    \begin{figure}[H]
                        \centering
                        \includegraphics[width=0.5\textwidth]{Solution/AMC_12_12B_2002_24_Solution_Figure_1.png}
                    \end{figure}
                    

 By the \href{https://artofproblemsolving.com/wiki/index.php?title=Pythagorean\_Theorem}{Pythagorean Theorem},  
 \begin{align*} AB = \sqrt{PA^2 + PB^2} & = \sqrt{24^2 + 32^2} = 40\\ BC = \sqrt{PB^2 + PC^2} & = \sqrt{32^2 + 28^2} = 4\sqrt{113}\\ CD = \sqrt{PC^2 + PD^2} & = \sqrt{28^2 + 45^2} = 53\\ DA = \sqrt{PD^2 + PA^2} & = \sqrt{45^2 + 24^2} = 51 \end{align*}


 The perimeter of $ABCD$ is $AB + BC + CD + DA = 4(36 + \sqrt{113}) \Rightarrow \mathrm{(E)}$.

 



\textbf{Solution 2}

Draw the diagram out. Notice the very peculiar sets of numbers $(24,32);(24,45);(28,45)$; these are simply multiples of the legs of well-known Pythagorean triples $(3,4,5);(8,15,17);(28,45,53)$, pointing to signs of possible right angles. We can assume that $\angle APB=\angle BPC=\angle CPD=\angle DPA=90^\circ$, so the area of the entire figure would be:

 
 \[A=\frac{1}{2}\cdot(24+28)(45+32)=\frac{1}{2}\cdot52\cdot77=2002\]

 Thus this is the correct case, so finding the side lengths of $ABCD$ by the Pythagorean Theorem yields $AB=40$, $BC=4\sqrt{113}$,$CD=53$,$DA=51$, so the perimeter is $40+4\sqrt{130}+53+51=144+4\sqrt{130}=\boxed{\textbf{(E) }4(36+\sqrt{113})}$.

 


	\item (\textit{AMC 12 12B 2002}) Let $f(x) = x^2 + 6x + 1$, and let $R$ denote the set of points $(x,y)$ in the coordinate plane such that 

 \[f(x) + f(y) \le 0 \qquad \text{and} \qquad f(x)-f(y) \le 0\]
The area of $R$ is closest to 


 $\mathrm{(A)}\ 21 \qquad\mathrm{(B)}\ 22 \qquad\mathrm{(C)}\ 23 \qquad\mathrm{(D)}\ 24 \qquad\mathrm{(E)}\ 25$

 \textbf{Answer: }E



\textbf{Solution 1}

The first condition gives us that
 \[x^2 + 6x + 1 + y^2 + 6y + 1 \le 0 \Longrightarrow (x+3)^2 + (y+3)^2 \le 16\]

 which is a \href{https://artofproblemsolving.com/wiki/index.php?title=Circle}{circle} centered at $(-3,-3)$ with \href{https://artofproblemsolving.com/wiki/index.php?title=Radius}{radius} $4$. The second condition gives us that 

 
 \[x^2 + 6x + 1 - y^2 - 6y - 1 \le 0 \Longrightarrow (x^2 - y^2) + 6(x-y) \le 0 \Longrightarrow (x-y)(x+y+6) \le 0\]

 Thus either

 
 \[x - y \ge 0,\quad x+y+6 \le 0\]

 or

 
 \[x - y \le 0,\quad x+y+6 \ge 0\]

                    \begin{figure}[H]
                        \centering
                        \includegraphics[width=0.5\textwidth]{Solution/AMC_12_12B_2002_25_Solution_Figure_1.png}
                    \end{figure}
                    Each of those lines passes through $(-3,-3)$ and has slope $\pm 1$, as shown above. Therefore, the area of $R$ is half of the area of the circle, which is $\frac{1}{2} (\pi \cdot 4^2) = 8\pi \approx \boxed{\textbf{(E) }25}$.~SHEN KISLAY KAI

 



\textbf{Solution 2}

Similar to Solution 1, we proceed to get the area of the circle satisfying $f(x)+f(y) \le 0$, or $16 \pi$.

 Since $f(x)-f(y) \le 0 \implies f(x) \le f(y)$, we have that by symmetry, if $(x,y)$ is in $R$, then $(y,x)$ is not, and vice versa. Therefore, the shaded part of the circle above the line $y=x$ has the same area as the unshaded part below $y=x$, and the unshaded part above $y=x$ has the same area as the shaded part below $y=x$. This means that exactly half the circle is shaded, allowing us to divide by two to get $\frac{16 \pi }{2} = 8\pi \approx \boxed{\textbf{(E) }25}$. ~samrocksnature + ddot1 +Shen kislay kai

 \textbf{Note} The equation $\frac{(x-h)^2}{a^2} - \frac{(y-k)^2}{b^2} = 0 \hspace{1mm} \text{or} \hspace{1mm} \frac{(y-k)^2}{a^2} - \frac{(x-h)^2}{b^2} = 0$ is the equation of a degenerate hyperbola.

 


\end{enumerate}